\section{Nghiệp vụ và quy trình}\label{sec:nghiep-vu}
\subsection{Giới thiệu bài toán}
Olist là doanh nghiệp kết nối người bán nhỏ với các sàn thương mại điện tử lớn ở Brazil. Bài tập lớn xây dựng kho dữ liệu cho chủ đề \textbf{quản lý đơn hàng} (Order Management) trên bộ dữ liệu công khai Brazilian E-Commerce \cite{olist_ecommerce}, tích hợp thêm bộ Marketing Funnel \cite{olist_marketing} để phân tích cả quá trình thu hút người bán. Hai bộ dữ liệu được nối qua mã người bán (\code{seller\_id}), nên kho trả lời được câu hỏi xuyên quy trình: \emph{người bán đến từ kênh marketing nào thì bán được nhiều và giao đúng hạn hơn}.

\todo{Viết 2--3 câu về mục tiêu và phạm vi bài toán theo cách của nhóm; thêm 1 đoạn giới thiệu Olist từ nguồn công khai (trích dẫn nguồn đã đọc).}

\paragraph{Các nhóm câu hỏi phân tích chính.}
\begin{itemize}
  \item \textbf{Bán hàng}: doanh thu (GMV), số đơn, giá trị đơn trung bình theo thời gian, danh mục, vùng; tăng trưởng và quy luật 80/20.
  \item \textbf{Xử lý đơn}: mỗi chặng (duyệt, bàn giao vận chuyển, giao) mất bao lâu, tỷ lệ giao trễ và hủy, vùng nào chậm.
  \item \textbf{Thanh toán}: cơ cấu hình thức thanh toán, số kỳ trả góp.
  \item \textbf{Trải nghiệm khách}: giao trễ ảnh hưởng điểm đánh giá thế nào, khách có quay lại không.
  \item \textbf{Người bán}: seller nào doanh thu cao nhưng hay trễ.
  \item \textbf{Thu hút người bán}: tỷ lệ chuyển đổi từ khách tiềm năng (MQL) sang seller theo kênh, thời gian chốt deal, thời gian đến đơn đầu tiên.
\end{itemize}

\paragraph{Mô hình kinh doanh (bản nháp).} Hình \ref{fig:bmc} tóm tắt mô hình kinh doanh của Olist theo 9 ô của Business Model Canvas, suy luận từ dữ liệu và mô tả công khai.
\hinh{3-2-3_business_model_canvas.png}{Business Model Canvas của Olist (bản nháp)}{fig:bmc}{0.98}
\todo{Nhóm kiểm tra và sửa Business Model Canvas, đặc biệt ô \emph{dòng doanh thu} và \emph{đối tác chính}: chỉ có suy luận, cần dẫn chứng từ nguồn công khai về Olist (dữ liệu chỉ có giá hàng và phí vận chuyển, không có doanh thu của Olist). Vẽ lại bằng draw.io nếu muốn.}

\subsection{Quy trình nghiệp vụ}
Chu trình xử lý đơn hàng gồm các bước: khách đặt hàng và thanh toán, sàn duyệt thanh toán, người bán chuẩn bị và bàn giao cho đơn vị vận chuyển, vận chuyển giao hàng, khách nhận hàng và đánh giá; đơn có thể bị hủy hoặc không khả dụng. Mỗi mốc là một cột thời gian trong dữ liệu nguồn (Hình \ref{fig:bpmn-don}); các ô vàng đánh dấu nơi \textbf{kho dữ liệu thu thập} để phân tích.
\hinh{3-2-1_bpmn_don_hang.png}{Quy trình xử lý đơn hàng và các điểm thu thập dữ liệu của kho (bản nháp)}{fig:bpmn-don}{1.0}
Quy trình thu hút người bán đi từ khách tiềm năng (MQL) để lại thông tin trên trang đích, đội sales (SDR, SR) tư vấn và chốt thương vụ, người bán ký hợp đồng và nhận đơn đầu tiên (Hình \ref{fig:bpmn-seller}). Trong dữ liệu: 8.000 MQL, 842 thương vụ được chốt (10,5\%), trong đó 379 seller có đơn bán được (45,0\% số đã chốt).
\hinh{3-2-2_bpmn_thu_hut_seller.png}{Quy trình thu hút người bán và các điểm thu thập dữ liệu của kho (bản nháp)}{fig:bpmn-seller}{1.0}
Các hoạt động của doanh nghiệp được tổ chức thành cây nghiệp vụ ở Hình \ref{fig:cay-nghiep-vu}.
\hinh{3-2-5_cay_nghiep_vu.png}{Cây nghiệp vụ của Olist (bản nháp)}{fig:cay-nghiep-vu}{0.8}
\todo{Nhóm đọc, sửa cho đúng cách hiểu nghiệp vụ và vẽ lại BPMN bằng draw.io (đúng ký hiệu BPMN), giữ nguyên các điểm thu thập dữ liệu. Nguồn nháp: huong-dan/05-nghiep-vu-nhap.md.}

\begin{table}[htbp]\centering\small\caption{Các mốc của quy trình đơn hàng và nơi kho lưu}\label{tab:moc-don}
\begin{tabular}{p{3.2cm}p{4.6cm}p{3.4cm}p{3.6cm}}\toprule
\textbf{Bước} & \textbf{Cột nguồn} & \textbf{Nơi kho lưu} & \textbf{Phân tích được}\\\midrule
Đặt hàng, thanh toán & \code{order\_purchase\_timestamp}, \code{payment\_type} & \code{purchase\_ts}, \code{fact\_payment} & số đơn, cơ cấu thanh toán\\
Duyệt & \code{order\_approved\_at} & \code{approval\_hours} & tốc độ duyệt\\
Bàn giao vận chuyển & \code{order\_delivered\_carrier\_date} & \code{handling\_days} & seller xử lý nhanh hay chậm\\
Giao khách & \code{order\_delivered\_customer\_date} & \code{transit\_days}, \code{lead\_days}, \code{delay\_days} & vận chuyển, đúng hạn\\
Đánh giá & \code{review\_score} & \code{fact\_review} & điểm trung bình, ảnh hưởng của trễ\\\bottomrule
\end{tabular}\end{table}

\subsection{Điểm phân tích trên quy trình}\label{sec:diem-phan-tich}
Trên hai quy trình (Hình \ref{fig:bpmn-don} và \ref{fig:bpmn-seller}), nhóm đánh dấu 11 điểm phân tích mã A1 đến A11. Mỗi điểm gắn một bước của quy trình với dashboard, KPI và khía cạnh phân tích sẽ dùng (Bảng \ref{tab:diem-phan-tich}); nhờ vậy dashboard nào cũng truy ngược được về một bước nghiệp vụ cụ thể.
\begin{table}[htbp]\centering\scriptsize\caption{Điểm phân tích trên quy trình: bước BPMN, dashboard, KPI và khía cạnh phân tích}\label{tab:diem-phan-tich}
\begin{tabular}{p{0.76cm}p{3.72cm}p{1.60cm}p{3.89cm}p{3.63cm}}\toprule
\textbf{Mã} & \textbf{Bước trên quy trình (BPMN)} & \textbf{Dashboard} & \textbf{KPI} & \textbf{Khía cạnh phân tích}\\\midrule
A1 & Quy trình đơn: đặt hàng và thanh toán & D1, D2, D5 & GMV, Orders, AOV, cơ cấu thanh toán, số kỳ trả góp & thời gian, sản phẩm, địa lý, hình thức thanh toán\\
A2 & Quy trình đơn: duyệt thanh toán & D3 & approval\_hours (trung bình 10,3 giờ) & thời gian, trạng thái đơn\\
A3 & Quy trình đơn: seller chuẩn bị và bàn giao & D3, D6 & handling\_days (trung bình 2,9 ngày) & seller, vùng, thời gian\\
A4 & Quy trình đơn: vận chuyển và giao hàng & D3, D4 & Avg lead days, On-time rate, Promised vs actual, Freight ratio & bang và vùng, khoảng cách, thời gian\\
A5 & Quy trình đơn: khách đánh giá & D5 & Avg review, \% negative, Repeat customer rate & mức trễ giao hàng, nhóm ngành, khách\\
A6 & Quy trình đơn: hủy hoặc không khả dụng & D1, D3 & Cancel rate, phễu trạng thái đơn & trạng thái, thời gian\\
A7 & Quy trình seller: khách tiềm năng để lại thông tin (MQL) & D6 & MQLs theo tháng liên hệ đầu & kênh (origin), thời gian\\
A8 & Quy trình seller: SDR sàng lọc, SR tư vấn & D6 & Avg days to close & nhân viên sales (SDR, SR), phân khúc\\
A9 & Quy trình seller: chốt thương vụ & D6 & Won deals, MQL to won rate & kênh, phân khúc kinh doanh\\
A10 & Quy trình seller: seller bán đơn đầu tiên & D6, D2 & Sellers with first sale, Won to first sale rate, GMV 90 ngày đầu & kênh, thời gian, seller\\
A11 & Vận hành dữ liệu: nạp và kiểm thử kho & D7 & số dòng nạp, thời gian chạy, số dòng bị loại, số kiểm thử đạt & lô nạp, bảng\\
\bottomrule\end{tabular}\end{table}


\subsection{Business Model Canvas và câu hỏi quản trị}\label{sec:bmc-cau-hoi}
Từ mỗi khối của Business Model Canvas (Hình \ref{fig:bmc}), nhóm rút ra một câu hỏi quản trị, chỉ số trả lời câu hỏi đó và dashboard hiển thị chỉ số (Bảng \ref{tab:bmc-cau-hoi}). Dữ liệu chỉ có giá hàng và phí vận chuyển nên hai ô \emph{dòng doanh thu} và \emph{cấu trúc chi phí} dùng GMV và tỷ lệ phí vận chuyển làm chỉ số thay thế; chi phí marketing và doanh thu thực của Olist không có trong dữ liệu.
\begin{table}[htbp]\centering\scriptsize\caption{Business Model Canvas: câu hỏi quản trị, chỉ số và dashboard tương ứng}\label{tab:bmc-cau-hoi}
\begin{tabular}{p{2.66cm}p{4.08cm}p{5.32cm}p{1.95cm}}\toprule
\textbf{Khối Canvas} & \textbf{Câu hỏi quản trị} & \textbf{Chỉ số} & \textbf{Dashboard}\\\midrule
Phân khúc khách hàng & Ai mua và mua ở đâu; ai là seller bán tốt? & số khách, Repeat customer rate, tỷ trọng vùng (Sudeste 65,4\% GMV) & D1, D4, D5\\
Giá trị đề xuất & Giao có đúng hạn và đủ rẻ để seller nhỏ vẫn bán được? & On-time rate, Promised vs actual, Freight ratio (16,6\%) & D3, D4\\
Kênh & Kênh marketing nào đưa về seller chốt được? & MQLs, MQL to won rate theo kênh & D6\\
Quan hệ khách hàng & Khách có hài lòng không; đội sales chốt nhanh không? & Avg review (4,09), \% negative, Avg days to close & D5, D6\\
Dòng doanh thu & Doanh số trên sàn tăng hay giảm? & GMV, Orders, AOV (dữ liệu không có doanh thu của Olist) & D1, D2\\
Nguồn lực chính & Nền tảng và đội sales đang phục vụ bao nhiêu seller? & số seller có đơn, số thương vụ trên mỗi nhân viên sales & D6\\
Hoạt động chính & Xử lý đơn và thu hút seller mất bao lâu ở từng chặng? & approval\_hours, handling\_days, Avg lead days, Avg days to close & D3, D6\\
Đối tác chính & Seller và đơn vị vận chuyển nào làm chậm đơn? & tỷ lệ trễ theo seller và bang, scorecard seller & D3, D4, D6\\
Cấu trúc chi phí & Chi phí vận chuyển chiếm bao nhiêu giá hàng? & Freight ratio theo vùng và nhóm ngành (dữ liệu không có chi phí marketing) & D2, D4\\
\bottomrule\end{tabular}\end{table}

\todo{Nhóm kiểm tra các câu hỏi quản trị theo cách hiểu nghiệp vụ của nhóm; nguồn nháp: huong-dan/05-nghiep-vu-nhap.md mục 8, 9.}
